% 所有版式使用完全相同的示例内容；不包含实验测量值。
\ReportDraftfalse
\ReportAppendixfalse
\newcommand{\ReportTitle}{中文、公式与图表排版最小示例}
\newcommand{\ReportTitleEN}{Minimal Example of Chinese Text, Equations and Figures}
\newcommand{\ReportAuthor}{\ReportBlank{2.2cm}}
\newcommand{\ReportAuthorEN}{\ReportBlank{2.2cm}}
\newcommand{\ReportStudentID}{\ReportBlank{2.5cm}}
\newcommand{\ReportTeacher}{\ReportBlank{2.2cm}}
\newcommand{\ReportExperimentDate}{\ReportBlank{2.5cm}}
\newcommand{\ReportAbstractCN}{本示例用于比较五种物理实验报告的排版风格，采用完全相同的中文段落、数学公式、符号表格、解析示意图及参考文献，观察题头、摘要区域、正文栏数和图表尺寸的区别。表格仅列出物理量符号及单位，曲线由余弦函数直接绘制，不包含实测数据。通过编译日志、中文显示、引用编号和页面边界检查，评估各版式的可读性。示例仅供样式选择，不能作为正式实验报告提交。}
\newcommand{\ReportAbstractEN}{This example compares five layouts for physics laboratory reports using identical Chinese text, equations, a symbol table, an analytical diagram and a reference. It illustrates differences in the title area, abstract arrangement, body columns and figure size. The table contains symbols and units; the curve is an analytical cosine function. No experimental measurements are included. The example is intended for choosing a layout and is not a report of an experiment.}
\newcommand{\ReportKeywordsCN}{中文排版；数学公式；图表；版式比较；字体}
\newcommand{\ReportKeywordsEN}{Chinese typesetting; equations; figures; layout comparison; fonts}
